Source : modèle communal du bloc (note8p/calculs/foncier_calculs.py) : DVF 2023-2024 (prix), REI 2025 (taxe foncière), Insee Analyses n° 91 (taxe / valeur 2017), comptes de patrimoine 2025 (part du sol).\par Taxe foncière bâtie des logements des particuliers rapportée à leur valeur de marché, par dixième de logements classés selon le prix au m² de leur commune (métropole hors Alsace-Moselle) ; abscisse : prix médian du dixième. Après réforme : base sol + bâti/4, part communale à recette constante, part nationale de 15 Md€ ; valeurs avant capitalisation. Même taux de la valeur pour tous les logements d'une commune.
